\vspace*{2.5cm}\noindent\begin{tabular}{@{}|p{1cm}p{11.9cm}|@{}} \hline
\multicolumn{2}{|l|}{Grille des points d'arrêt du {\normalsize \robot}~:}  \\ \hline \hline
\idt   dg    & \apxFontSize \underline{d}essine la \underline{g}rille des points d'arrêt du {\normalsize \rm\robot}\ (peut être interrompue par
%touche Esc (touche d'echappement)
\begin{picture}(14,12)(0, 2) 
\put(0,0){\framebox(14,12)[lb]{}}
\put(1.5,2){\makebox(0,0)[lb]{\tiny Esc}}
\end{picture}\hspace{0.3ex},
%touche §
\begin{picture}(12,12)(0, 2)
\put(0,0){\framebox(12,12)[lb]{}}
\put(1.5,2){\makebox(0,0)[lb]{\tiny \S}}
\end{picture}\hspace{0.3ex}
ou
%touche #
\begin{picture}(12,12)(0, 2)
\put(0,0){\framebox(12,12)[lb]{}}
\put(1.5,2){\makebox(0,0)[lb]{\tiny \#}}
\end{picture}\hspace{0.3ex}). \\
\idt   lg$(n)$ & \apxFontSize définit  la  nouvelle \underline{l}ongueur $n$ des côtés d'un carreau de la 
          \underline{g}rille ($(12.0\:\times\:12.0)$ pixels par défaut, $n\:\geq \: 0.0$). \\
\idt   lgX$(n)$& \apxFontSize définit  la  nouvelle \underline{l}ongueur $n$ du côté horizontal (\underline{X}) 
           d'un carreau de la \underline{g}rille ($12.0$ pixels par défaut, $n\:\geq \: 0.0$).\\
\idt   lgY$(n)$& \apxFontSize définit la nouvelle \underline{l}ongueur n du côté vertical (\underline{Y}) d'un carreau 
         de la \underline{g}rille ($12.0$ pixels par défaut, $n\:\geq \: 0.0$). \\
\idt   lgRT$(r, t)$ & \apxFontSize \hspace{1em} définit les nouvelles \underline{l}ongueurs des côtés d'un carreau 
         de la \underline{g}rille par les coordonnées polaires \underline{$r$} et \underline{$t$} (la longueur 
		 $r\:\geq\:0.0$ de la diagonale du carreau est exprimée en pixels et l'angle  $t$ que fait cette diagonale avec le 
		 côté horizontal est exprimé en degrés $(0.0\:\leq\:t\leq\:90.0)$ ).
\\ \hline
\idt   valLgX& \apxFontSize est égal à la \underline{val}eur (de type {\idt real}) de la \underline{l}on\underline{g}ueur 
du côté horizontal (\underline{X}) d'un carreau de la grille (comptée en pixels). \\
\idt   valLgY& \apxFontSize est égal à la \underline{val}eur (de type {\idt real}) de la \underline{l}on\underline{g}ueur 
du côté vertical (\underline{Y}) d'un carreau de la grille (comptée en pixels). \\
\idt   valLgR& \apxFontSize est égal à la \underline{val}eur (de type {\idt real}) de la \underline{l}on\underline{g}ueur 
de la diagonale (\underline{R}) d'un carreau de la grille (comptée en pixels). \\
\idt   valLgT& \apxFontSize est égal à la \underline{val}eur (de type {\idt real}) de \underline{l}'an\underline{g}le (\underline{T}) de la diagonale 
                       d'un carreau de la grille avec son côté horizontal (comptée en degrés).  \\
\hline \end{tabular}

\noindent\begin{tabular}{@{}|p{1cm}p{11.9cm}|@{}} \hline
\idt   valX  & \apxFontSize est égal à la \underline{val}eur (de type {\idt real}) de l'abscisse absolue (\underline{X}) du 
                           {\normalsize \rm\robot}\ (comptée en pixels). \\
\idt   valY  & \apxFontSize est égal à la \underline{val}eur (de type {\idt real}) de l'ordonnée absolue (\underline{Y}) du 
           {\normalsize \rm\robot}  (comptée en pixels), de  sorte  que 
           {\idt X $:=$ valX$;$ Y := valY$;$} et plus tard {\idt lg$($1.0$);$ deplacer$($X$,$ Y$);$}   replace le 
           {\normalsize \rm\robot}\ à sa place initiale sur l'écran). \\ \hline
\end{tabular}

\noindent\begin{tabular}{@{}|p{1cm}p{11.9cm}|@{}} \hline
\multicolumn{2}{|l|}{Pivotements relatifs sur place~:}  \\ \hline \hline
 \idt     phtd  & \apxFontSize  fait \underline{p}ivoter le {\normalsize \rm\robot}\ sur place d'un \underline{h}uitième 
                          de \underline{t}our à \underline{d}roite. \\
\idt      pqtd  & \apxFontSize  fait \underline{p}ivoter le {\normalsize \rm\robot}\  sur place d'un \underline{q}uart de 
                          \underline{t}our à \underline{d}roite. \\
 \idt     p3htd & \apxFontSize  fait \underline{p}ivoter le {\normalsize \rm\robot}\  sur place de \underline{3} \underline{h}uitièmes de 
                             \underline{t}our à \underline{d}roite. \\
\idt      phtg  & \apxFontSize  fait \underline{p}ivoter le {\normalsize \rm\robot}\  sur place d'un \underline{h}uitième de \underline{t}our 
                           à \underline{g}auche.  \\
\idt      pqtg  & \apxFontSize fait \underline{p}ivoter le {\normalsize \rm\robot}\  sur place d'un \underline{q}uart de \underline{t}our à 
                             \underline{g}auche. \\
\idt      p3htg & \apxFontSize  fait \underline{p}ivoter le {\normalsize \rm\robot}\  sur place de \underline{3} \underline{h}uitièmes de 
                           \underline{t}our à \underline{g}auche. \\
\idt      pdt   & \apxFontSize  fait \underline{p}ivoter le {\normalsize \rm\robot}\  sur place d'un \underline{d}emi-\underline{t}our.  \\ \hline
\end{tabular}

\noindent\begin{tabular}{@{}|p{1cm}p{11.9cm}|@{}} \hline
\multicolumn{2}{|l|}{Pivotements absolus sur place~: }  \\ \hline \hline
 \idt  pn    & \apxFontSize   fait \underline{p}ivoter le {\normalsize \rm\robot}\  sur place vers le nord (\underline{N}).   \\                   
 \idt  pne   & \apxFontSize   fait \underline{p}ivoter le {\normalsize \rm\robot}\ sur place vers le nord-est (\underline{NE}).  \\                
\idt   pe    & \apxFontSize   fait \underline{p}ivoter le {\normalsize \rm\robot}\ sur place vers l'est (\underline{E}). \\                
\idt   pse   & \apxFontSize   fait \underline{p}ivoter le {\normalsize \rm\robot}\ sur place vers le sud-est (\underline{SE}). \\      
\idt   ps    & \apxFontSize   fait \underline{p}ivoter le {\normalsize \rm\robot}\ sur place vers le sud (\underline{S}). \\           
\idt   psw   & \apxFontSize   fait \underline{p}ivoter le {\normalsize \rm\robot}\ sur place vers le sud-ouest (\underline{SW}).  \\ 
\idt   pw    & \apxFontSize   fait \underline{p}ivoter le {\normalsize \rm\robot}\ sur place vers l'ouest (\underline{W}).   \\ 
\idt   pnw   & \apxFontSize   fait \underline{p}ivoter le {\normalsize \rm\robot}\ sur place vers le nord-ouest (\underline{NW}).  \\ 
\idt     pivoter$(o)$ & \apxFontSize \hspace{2em}  fait \underline{pivoter} le {\normalsize \rm\robot}\ dans 
                 l'\underline{o}rientation  indiquée 
                    par la valeur de l'entier {\idt  $o$}  modulo 8 égale
                      à  {\idt nord} $=$ 0,  {\idt nordEst} $=$ 1,  {\idt est} $=$ 2, {\idt sudEst} $=$ 3, {\idt sud} $=$ 4, 
					  {\idt sudOuest} $=$ 5, {\idt ouest} $=$ 6 ou {\idt nordOuest} $=$ 7. \\                     
\idt     rt$(t)$ & \apxFontSize   change le repère et l'orientation du {\normalsize \rm\robot}\ pour qu'il effectue une \underline{r}o\underline{t}ation sur place d'un angle de
           $t$ degrés (dans le sens trigonométrique si $t > 0$  et  dans le sens des aiguilles 
           d'une montre si $t < 0$).   \\ \hline                                                
\idt     valOrientation & \apxFontSize \hspace{4em} est égal  à la \underline{val}eur (de type {\idt integer}) 
                         de l'\underline{orientation} du {\normalsize \rm\robot}\ codée par
                      {\idt nord} $=$ 0,  {\idt nordEst} $=$ 1,  {\idt est} $=$ 2, {\idt sudEst} $=$ 3, {\idt sud} $=$ 4, 
					  {\idt sudOuest} $=$ 5, {\idt ouest} $=$ 6 et {\idt nordOuest} $=$ 7.  Les directions {\idt nord}, 
					  {\idt est}, {\idt sud} et {\idt ouest}
                       sont parallèles aux bords de l'écran, le {\idt nord} en haut. Les directions {\idt nordEst}, 
					  {\idt sudEst}, {\idt sudOuest} et {\idt nordOuest}
                       correspondent à l'orientation des diagonales d'un carreau de la grille.\\ \hline             
\end{tabular}

\noindent\begin{tabular}{@{}|p{1cm}p{11.9cm}|@{}} \hline
\multicolumn{2}{|l|}{Crayon~: }  \\ \hline \hline
 \idt    bc    & \apxFontSize    \underline{b}aisse le \underline{c}rayon pour dessiner lors des déplacements.  \\ 
\idt    lc    & \apxFontSize    \underline{l}ève le \underline{c}rayon pour éviter de dessiner lors des déplacements.  \\ 
\idt    dc    & \apxFontSize    \underline{d}essine une petite \underline{c}roix (de l'épaisseur du crayon), sans déplacer le {\normalsize \rm\robot}.  \\ 
\idt    dp    & \apxFontSize    \underline{d}essine un \underline{p}oint (de l'épaisseur du crayon), sans déplacer le {\normalsize \rm\robot} . \\      
\idt    ec$(n)$ & \apxFontSize   définit la nouvelle \underline{é}paisseur $n$ du \underline{c}rayon (1 pixel par  défaut).  \\ \hline     
\idt   valEc & \apxFontSize    est égal à la \underline{val}eur (de type {\idt integer}) de l'\underline{é}paisseur du \underline{c}rayon. \\                  
\idt    valBc & \apxFontSize   est égal à la \underline{val}eur booléenne {\idt true} (vrai) si et seulement si le \underline{c}rayon est \underline{b}aissé,  
           sinon à {\idt false} (faux). \\ \hline                                               
\end{tabular}

\noindent\begin{tabular}{@{}|p{1cm}p{11.9cm}|@{}} \hline
\multicolumn{2}{|l|}{Déplacements (avec ou sans tracé selon que le crayon est levé ou baissé)~: }  \\ \hline \hline
 \idt  av    & \apxFontSize    fait \underline{av}ancer le {\normalsize \rm\robot}\ tout droit jusqu'au prochain point d'arrêt sur la grille. \\ 
 \idt     vd    & \apxFontSize    fait tourner le {\normalsize \rm\robot}\  par un \underline{v}irage d'un quart de cercle à \underline{d}roite.  \\ 
 \idt     vg    & \apxFontSize    fait tourner le {\normalsize \rm\robot}\  par un \underline{v}irage d'un quart de cercle à \underline{g}auche. \\ 
  \idt    avf$(n)$ & \apxFontSize    fait \underline{av}ancer le {\normalsize \rm\robot}\  tout droit $n$ \underline{f}ois 
           (sans rien faire si $n = 0$ et en reculant si $n < 0$). \\ 
 \idt     tl$(l)$ & \apxFontSize    change le repère (mais pas l'angle des directions NE-E)  et fait avancer le {\normalsize \rm\robot}\  par une
            \underline{t}rans\underline{l}ation  de $l$ pixels dans la direction où il se trouve (si $l > 0$) ou 
           dans la direction opposée (si $l < 0$). \\ 
 \idt     tracer$(X, Y)$ & \apxFontSize   \hspace{3em} \underline{trace} un segment de droite (si le crayon est baissé) de la  position 
           courante du {\normalsize \rm\robot}\  au point de coordonnées $(X, Y)$  et déplace le {\normalsize \rm\robot}\  en 
           ce point.\\ \hline
 \idt     valLgTl & \apxFontSize \hspace{1em}    est égal à la \underline{val}eur (de type {\idt real}) de la \underline{l}on\underline{g}ueur du segment tracé
           par la dernière \underline{t}rans\underline{l}ation par {\idt av}  ou {\idt tl} (comptée en pixels).   \\ \hline      
\end{tabular}

\noindent\begin{tabular}{@{}|p{1cm}p{11.9cm}|@{}} \hline
\multicolumn{2}{|l|}{Positionnement au centre et sur les bords du cadre~: }  \\ \hline \hline
\idt       ce    & \apxFontSize  place le {\normalsize \rm\robot}\  et {\em  le centre de la grille} au \underline{ce}ntre de {\em l'écran}, sans changer l'orientation du 
                       {\normalsize \rm\robot}\ et sans rien dessiner.  \\ 
\idt       eh    & \apxFontSize  place  le  {\normalsize \rm\robot}\  sur la grille \underline{e}n \underline{h}aut,  sans  changer  sa position verticale, ni son 
           orientation et sans rien dessiner.  \\ 
 \idt      eb    & \apxFontSize  place  le  {\normalsize \rm\robot}\  sur la grille  \underline{e}n \underline{b}as,  sans  changer  sa  position verticale, ni son 
           orientation et sans rien dessiner.  \\                 
 \idt      ag    & \apxFontSize  place  le  {\normalsize \rm\robot}\  sur la grille  \underline{à} \underline{g}auche, sans changer sa position horizontale, ni son 
           orientation et sans rien dessiner.  \\ 
 \idt      ad    & \apxFontSize  place  le  {\normalsize \rm\robot}\ sur la grille   \underline{à} \underline{d}roite, sans changer sa position horizontale, 
          ni son orientation et sans rien dessiner.  \\ \hline
 \end{tabular}

\noindent\begin{tabular}{@{}|p{1cm}p{11.9cm}|@{}} \hline
\multicolumn{2}{|l|}{Repère cartésien~: }  \\ \hline \hline
 \idt       translater$(X, Y) $& \apxFontSize   \hspace{5em} \underline{translate}  le  {\normalsize \rm\robot}\   de  $X$ carreaux de la grille 
 ho\-ri\-zon\-ta\-le\-ment (à droite si $X > 0$
         et à gauche si $X < 0$) et
          $Y$ carreaux de la grille verticalement (en haut si $Y > 0$ et en bas si $Y < 0$), sans changer son orientation et  sans  rien 
           dessiner.   \\
\idt         deplacer$(X, Y)$ & \apxFontSize  \hspace{4em} \underline{déplace}  le {\normalsize \rm\robot}\  au point de coordonnées  (X, Y) sur la grille, 
           sans changer son orientation ni rien dessiner, l'origine du repère cartésien étant placée au 
           centre de la grille. \\ \hline                                          
\end{tabular}

\noindent\begin{tabular}{@{}|p{1cm}p{11.9cm}|@{}} \hline
\multicolumn{2}{|l|}{Test booléen de l'orientation du {\normalsize \rm\robot}~:  }  \\ \hline \hline
 \idt    tn    &   \apxFontSize  \underline{t}este si le {\normalsize \rm\robot}\  est orienté au nord (\underline{N}). \\ 
\idt    tne   &   \apxFontSize  \underline{t}este si le {\normalsize \rm\robot}\  est orienté au nord-est (\underline{NE}). \\ 
\idt    te    &   \apxFontSize  \underline{t}este si le {\normalsize \rm\robot}\  est orienté à l'est (\underline{E}).\\ 
\idt    tse   &   \apxFontSize  \underline{t}este si le {\normalsize \rm\robot}\  est orienté au sud-est (\underline{SE}).\\ 
\idt    ts    &   \apxFontSize  \underline{t}este si le {\normalsize \rm\robot}\  est orienté au sud (\underline{S}).\\ 
 \idt   tsw   &   \apxFontSize  \underline{t}este si le {\normalsize \rm\robot}\  est orienté au sud-ouest (\underline{SW}).\\ 
\idt    tw    &   \apxFontSize  \underline{t}este si le {\normalsize \rm\robot}\  est orienté à l'ouest (\underline{W}).\\ 
\idt    tnw   &   \apxFontSize  \underline{t}este si le {\normalsize \rm\robot}\  est orienté au nord-ouest (\underline{NW}).  \\ \hline                                    
\end{tabular}

\noindent\begin{tabular}{@{}|p{1cm}p{11.9cm}|@{}} \hline
\multicolumn{2}{|l|}{Test booléen de la position du {\normalsize \rm\robot}~:  }  \\ \hline \hline
\idt         tb &   \apxFontSize    \underline{t}este si le {\normalsize \rm\robot}\  est au \underline{b}ord du cadre (un déplacement du {\normalsize \rm\robot}\  en avant 
           par {\idt av} le faisant sortir du cadre).   \\ \hline                           
\end{tabular}

\noindent\begin{tabular}{@{}|p{1cm}p{11.9cm}|@{}} \hline
\multicolumn{2}{|l|}{Temps~:  }  \\ \hline \hline
\idt   heure$($H$,$~Mi$,$~S$)$  & \apxFontSize \hspace{4em}  affecte aux  variables  {\idt H},  {\idt Mi}  et  {\idt S}  trois  entiers qui sont 
           l'heure (entre 0 et 23),  la minute (entre 0 et 59) et la seconde (entre 
           0 et 59) présentes.   \\
\idt   date$($J$,$~Mo$,$~A$)$  & \apxFontSize \hspace{4em}  affecte aux variables {\idt J}, {\idt Mo} et {\idt A} trois entiers qui sont le jour 
           (entre 1 et 31), le mois (entre 1 et 12) et l'année courante. \\ \hline
\idt   delai$(t)$ & \apxFontSize \hspace{1ex} attendre pendant un \underline{délai} de $t$ secondes ($t$ de type {\idt real}).\\ \hline                                        
\end{tabular}

\noindent\begin{tabular}{@{}|p{1cm}p{11.9cm}|@{}} \hline
\multicolumn{2}{|l|}{Interruption de l'exécution du programme~:  }  \\ \hline \hline
\mbox{%touche Esc (touche d'echappement)
\begin{picture}(14,12)(0, 2) 
\put(0,0){\framebox(14,12)[lb]{}}
\put(1.5,2){\makebox(0,0)[lb]{\tiny Esc}}
\end{picture} ou 
%touche §
\begin{picture}(12,12)(0, 2)
\put(0,0){\framebox(12,12)[lb]{}}
\put(1.5,2){\makebox(0,0)[lb]{\tiny \S}}
\end{picture} ou 
%touche #
\begin{picture}(12,12)(0, 2)
\put(0,0){\framebox(12,12)[lb]{}}
\put(1.5,2){\makebox(0,0)[lb]{\tiny \#}}
\end{picture}\hspace{0.3ex}}
 & \apxFontSize \hspace{5em} stoppe le {\normalsize \rm\robot} au cours de l'exécution de l'une de ses commandes. \\ \hline           
\end{tabular}

\noindent\begin{tabular}{@{}|p{1cm}p{11.9cm}|@{}} \hline
\multicolumn{2}{|l|}{ Mise au point du programme~:  }  \\ \hline \hline
\idt            st  &   \apxFontSize \underline{st}oppe le {\normalsize \rm\robot} qui s'arrête définitivement en tapant l'une des touches
\mbox{%touche Esc (touche d'echappement)
\begin{picture}(14,12)(0, 2) 
\put(0,0){\framebox(14,12)[lb]{}}
\put(1.5,2){\makebox(0,0)[lb]{\tiny Esc}}
\end{picture}\hspace{0.3ex},}
%touche §
\begin{picture}(12,12)(0, 2)
\put(0,0){\framebox(12,12)[lb]{}}
\put(1.5,2){\makebox(0,0)[lb]{\tiny \S}}
\end{picture}\hspace{0.3ex}
ou
%touche #
\begin{picture}(12,12)(0, 2)
\put(0,0){\framebox(12,12)[lb]{}}
\put(1.5,2){\makebox(0,0)[lb]{\tiny \#}}
\end{picture}\hspace{0.3ex})
 ou redémarre en enfonçant une autre touche du clavier ou le bouton de la souris. Avant de redémarrer,
 l'utilisateur a la possibilité de diminuer (avec la touche
\mbox{\begin{picture}(12,12)(0, 2)
\put(0,0){\framebox(12,12)[lb]{}}
\put(1.5,2){\makebox(0,0)[lb]{\tiny $-$}}
\end{picture}\hspace{0.3ex})} ou d'augmenter (avec la touche
\mbox{\begin{picture}(12,12)(0, 2)
\put(0,0){\framebox(12,12)[lb]{}}
\put(1.5,2){\makebox(0,0)[lb]{\tiny $+$}}
\end{picture}\hspace{0.3ex})}
la vitesse du {\normalsize \rm\robot} (entre 0 et 10, la vitesse 0 faisant passer en mode pas à pas).\\ 
\idt            pp   &   \apxFontSize exécute la suite du programme en mode \underline{p}as à \underline{p}as, à la vitesse 0. \\ 
\idt            ex   &   \apxFontSize exécute la suite du programme en mode \underline{ex}écution normale, à la vitesse 10. \\ 
\idt            vt$(v)$&   \apxFontSize exécute  la  suite  du  programme à la \underline{v}i\underline{t}esse $v$ (en pas à pas si 
$v \leq 0$,  avec un maximum de 10 si $v > 10$).  \\ 
\idt           montrerRobot  &   \apxFontSize \hspace{3em}  rend le {\normalsize \rm\robot}\  visible sur l'écran 
                 pendant le dessin (option par défaut).  \\ 
\idt           cacherRobot &   \apxFontSize \hspace{3em} rend le {\normalsize \rm\robot}\  invisible sur l'écran pendant le dessin.  \\ 
\idt           InterdireSorties &  \apxFontSize \hspace{5em}  \underline{interdit} les  \underline{sorties}  du {\normalsize \rm\robot}\ en  
                dehors  du cadre dans la suite du programme (option par défaut).  \\  
 \idt          autoriserSorties &   \apxFontSize \hspace{5em} \underline{autorise} les  \underline{sorties}  du {\normalsize \rm\robot}\ en  
                  dehors  du cadre dans la suite de l'exécution du programme. \\  \hline          
\idt           valVt &   \apxFontSize est égal à la \underline{val}eur (de type {\idt integer}) de la
\underline{v}i\underline{t}esse du {\normalsize \rm\robot}\ (entre 0 (mode pas à pas) et 10). \\ 
\idt           valVisible &   \apxFontSize \hspace{2em} est égal à la valeur booléenne {\idt true},  (respectivement 
                {\idt false}) si le {\normalsize \rm\robot}\  
          est \underline{visible} (respectivement invisible) sur l'écran. \\ 
\idt           valAutoriserSorties &   \apxFontSize \hspace{7em} est égal à la valeur booléenne 
{\idt true} (respectivement 
                {\idt false}) si 
               les \underline{sorties} en dehors du cadre sont \underline{autorisées} (respectivement interdites).    \\ \hline                                      
\end{tabular}

\noindent\begin{tabular}{@{}|p{1cm}p{11.9cm}|@{}} \hline
%
Couleur~: & \\ \hline \hline
%
\multicolumn{2}{|p{12.9cm}|}{Les couleurs sont  {\idt blanc} $= 0$,  {\idt jaune} $= 1$,  {\idt bleu} $= 2$,  {\idt violet} $= 3$, 
 {\idt vert} $= 4$, {\idt rouge} $= 5$, 
{\idt indigo} $= 6$, {\idt noir} $= 7$ (dans 
 l'ordre du plus clair au plus foncé sur un écran ayant au moins 8 niveaux de gris)
  et {\idt inverse = 8}.
 Sur un écran noir et blanc, et selon le modèle de l'ordinateur, toutes les couleurs  (sauf  blanc) 
 apparaissent en noir ou les couleurs sont représentées par des motifs répétés.}\\\hline  \hline
%
\idt      cc$(c)$ & \apxFontSize  définit la \underline{c}ouleur $c$ du \underline{c}rayon. Si la couleur est 
             {\idt inverse}, les tracés se feront en vidéo inverse (par exemple noir sur fond blanc et 
			 blanc sur fond noir).  \\
\hline
\end{tabular}

\noindent\begin{tabular}{@{}|p{1cm}p{11.9cm}|@{}} \hline
\idt      peindre & \apxFontSize  \hspace{1ex} \underline{peint} de la couleur du crayon la région blanche de surface 
            maximale où se trouve le {\normalsize \rm\robot}\ (la peinture peut être interrompue par
%touche Esc (touche d'echappement)
\begin{picture}(14,12)(0, 2) 
\put(0,0){\framebox(14,12)[lb]{}}
\put(1.5,2){\makebox(0,0)[lb]{\tiny Esc}}
\end{picture}\hspace{0.3ex},
%touche §
\begin{picture}(12,12)(0, 2)
\put(0,0){\framebox(12,12)[lb]{}}
\put(1.5,2){\makebox(0,0)[lb]{\tiny \S}}
\end{picture}\hspace{0.3ex}
ou
%touche #
\begin{picture}(12,12)(0, 2)
\put(0,0){\framebox(12,12)[lb]{}}
\put(1.5,2){\makebox(0,0)[lb]{\tiny \#}}
\end{picture}\hspace{0.3ex}, le \robot\ continuant ensuite son travail normalement). \\ 
\idt      cf$(c)$ & \apxFontSize  efface complètement le dessin puis remplit l'écran de la \underline{c}ouleur de 
            \underline{f}ond $c$, {\idt inverse} n'étant pas autorisé. Le {\normalsize \rm\robot}\ est placé au centre 
			de l'écran dans les conditions initiales (grille $(12 \times 12)$, visible, tourné vers le
			nord, sorties hors du cadre interdites, crayon noir de taille 1 baissé et exécution normale à
			la vitesse 10).    \\  \hline
\idt      valCf & \apxFontSize  est égal à la \underline{val}eur (de type {\idt integer}) de la \underline{c}ouleur du \underline{f}ond.  \\
\idt      valCc & \apxFontSize  est égal à la \underline{val}eur  (de type {\idt integer}) de la \underline{c}ouleur du \underline{c}rayon.  \\
\idt      ecranNoirEtBlanc & \apxFontSize \hspace{5em} est égal à {\idt true} si et seulement si l'écran est noir et 
           blanc (sinon {\idt false}). \\ \hline           
\end{tabular}

\noindent\begin{tabular}{@{}|p{1cm}p{11.9cm}|@{}} \hline
\multicolumn{2}{|l|}{Écriture dans la fenêtre de dessin du {\normalsize \rm\robot}~:  }  \\ \hline \hline
\idt    ecrire$(c)$ & \apxFontSize \hspace{1em} écrit la chaîne de caractères $c$ (entre 
           apostrophes {\cha $\scriptstyle '$} et {\cha $\scriptstyle '$}, où toute apostrophe {\cha $\scriptstyle '$} 
		   doit être doublée) à droite du {\normalsize \rm\robot}\ sans changer sa position ni celle de son crayon. \\
\idt    entierEnChaine$(i)$ & \apxFontSize \hspace{5em}   est la chaîne de caractères représentant l'entier $i$ en base 
     10. \\
 \idt   reelEnChaine$(r)$ & \apxFontSize \hspace{4.5em}   est  la  chaîne de caractères représentant le réel $r$ 
                    sous la forme {\idt $d.dddd$E$\pm ee$} où $d$ est un digit de la mantisse, $e$ est un digit de 
					l'exposant et {\idt E} se lit ``fois 10 à la puissance''. \\
\idt    booleenEnChaine$(b)$ & \apxFontSize \hspace{6.5em}  est  la  chaîne de caractères 
({\cha $\scriptstyle '$Vrai$\scriptstyle '$} ou {\cha $\scriptstyle '$Faux$\scriptstyle '$}) 
re\-pré\-sen\-tant le 
booléen $b$.\\
 \idt   concat$(c_1, c_2, \ldots, c_n)$ & \apxFontSize \hspace{6.5em}  est la chaîne de caractères  obtenue  en  
     \underline{concat}énant  les chaînes $c_1$, $c_2$, \ldots, $c_n$ les unes derrière les autres. \\ \hline
\idt    policeTexte$(p)$ & \apxFontSize \hspace{4em}  choisit la \underline{police} $p$ du \underline{texte} 
                                        (qui dépend du type d'ordinateur utilisé).  \\
 \idt   tailleTexte$(t)$ & \apxFontSize \hspace{4em}  choisit la \underline{taille} $t$ du \underline{texte}
                                          (qui dépend du type d'ordinateur utilisé).  \\
 \idt    styleTexte$(s)$ & \apxFontSize \hspace{4em}  choisit le \underline{style} du \underline{texte}  
                                             (qui dépend du type d'ordinateur utilisé).  \\   \hline                                                       
\idt    valPoliceTexte & \apxFontSize \hspace{4em}  \underline{val}eur de la \underline{police} courante du \underline{texte}.  \\
\idt   valTailleTexte & \apxFontSize \hspace{4em}  \underline{val}eur de la \underline{taille} courante du \underline{texte}. \\
\idt    valStyleTexte & \apxFontSize \hspace{4em}  \underline{val}eur du \underline{style} courant du \underline{texte}.  \\ \hline
\end{tabular}
\newpage
\noindent\begin{tabular}{@{}|p{1cm}p{11.9cm}|@{}} \hline
\multicolumn{2}{|l|}{Messages affichés en dessous de la fenêtre de déplacement du {\normalsize \rm\robot}~:  }  \\ \hline \hline
 \idt   message$(c)$ & \apxFontSize  \hspace{2em} efface le \underline{message} dans le cadre prévu en dessous de la 
           fenêtre de déplacement du {\normalsize \rm\robot}\ puis y écrit la chaîne de caractères $c$. \\
 \idt    marquerUnePause & \apxFontSize  \hspace{5em}  \underline{marque} \underline{une} \underline{pause} (en faisant clignoter le symbole $\hookleftarrow$
           dans la fenêtre des messages) en attendant que 
           le bouton de la souris ou qu'une touche du clavier soit enfoncé.\\
 \idt   effacerMessage & \apxFontSize  \hspace{4em}  \underline{efface} le \underline{message} dans le cadre prévu en dessous de la fenêtre 
           de déplacement du {\normalsize \rm\robot}. \\ \hline
\idt   lireCar$($C$)$ & \apxFontSize  \hspace{2em}  attend  (en faisant clignoter le symbole $\hookleftarrow$
           dans la fenêtre des messages) qu'une touche du clavier
             ou que le bouton de la souris soit enfoncé,  puis   \underline{lit} le  \underline{car}actère    correspondant  à la touche 
           enfoncée (ou le caractère nul {\idt chr$($0$)$} si le bouton de la souris a été
           enfoncé) et l'affecte  à la variable {\idt C} (de type caractère {\idt char}). \\ \hline  
 \idt  mementoCommandesClavier & \apxFontSize  \hspace{9em}  retourne  une  chaîne  de  caractères  rappelant  les 
           commandes de pilotage interactif du {\normalsize \rm\robot}. \\                            
\idt  valMessage & \apxFontSize  \hspace{2em} est égal à \underline{val}eur  (de type {\idt string}) du \underline{message} (de type {\idt string})  inscrit dans le cadre sous la 
           fenêtre de déplacement du {\normalsize \rm\robot}\ ou la chaîne vide {\cha $\scriptstyle '\scriptstyle '$} s'il n'y a aucun message.  \\ \hline           
\end{tabular}
\vfil
\noindent\begin{tabular}{@{}|p{1cm}p{11.9cm}|@{}} \hline
\multicolumn{2}{|l|}{Lecture et écriture dans la fenêtre texte de \pascal~:  }  \\ \hline \hline
\idt   modeTexte & \apxFontSize \hspace{2em} fait apparaître la fenêtre ``Texte"  de  \pascal\  (pour  réaliser  les 
          entrées-sorties avec les instructions {\idt write}, {\idt writeln}, {\idt read} et {\idt readln}). La réapparition de
          la fenêtre de dessin du {\normalsize \rm\robot}\  est automatique à la prochaine utilisation d'une 
          commande du {\normalsize \rm\robot}. \\ \hline                                                       
\end{tabular}
\vfil
\noindent\begin{tabular}{@{}|p{1cm}p{11.9cm}|@{}} \hline
\multicolumn{2}{|l|}{Impression et conservation du dessin du {\normalsize \rm\robot}~:  }  \\ \hline \hline
\idt    definirFormatImpression$(f)$  & \apxFontSize \hspace{10em}  \underline{définit}  le  
         \underline{format} $f =$ {\idt reduction} $= 0$, {\idt standard} $= 1$ ou {\idt agrandissement} $= 2$ de la
		 prochaine \underline{impression}. \\              
\idt    definirTitreImpression$(t)$  & \apxFontSize \hspace{9em}  \underline{définit} 
         la chaîne de caractères $t$ (entre apostrophes {\cha $\scriptstyle '$} et {\cha $\scriptstyle '$}, où toute apostrophe {\cha $\scriptstyle '$} 
		   doit être doublée) comme le \underline{titre} pour la prochaine \underline{impression}. \\     
\idt     imprimer  & \apxFontSize \hspace{1em}  \underline{imprime} le dessin dans le format et avec 
          le titre  (en  dessous du dessin)  tels qu'ils ont été dernièrement définis 
		  ({\idt standard} et  {\cha $\scriptstyle '$}{\cha $\scriptstyle '$} par défaut).  \\ \hline     
\idt    valFormatImpression  & \apxFontSize \hspace{7em}  est égal à la \underline{val}eur  (de type {\idt integer}) du 
             \underline{format} d'\underline{impression}.     \\                        
\idt    valTitreImpression  & \apxFontSize \hspace{7em}  est égal à la \underline{val}eur  (de type {\idt string}) du 
\underline{titre} d'\underline{impression}.  \\ \hline                            
\idt   copier  & \apxFontSize \underline{copie} le dessin sur disque (dépend de l'ordinateur utilisé).   \\ \hline                                                              
\end{tabular}
\newpage
\noindent\begin{tabular}{@{}|p{1cm}p{11.9cm}|@{}} \hline
\multicolumn{2}{|l|}{Sons~:  }  \\ \hline \hline
\idt bip   & \apxFontSize  émet un son bref. \\[-1mm]
\idt biiip & \apxFontSize  émet un son long. \\[-1mm]
\idt son$(f, t)$ & \apxFontSize \hspace{0.75ex} émet un \underline{son}
de fréquence $f$ {\scriptsize(comprise entre 15 et 15000)} pendant $t$ se\-con\-des
{\scriptsize($f$ et $t$ de type {\idt real}, peut être interrompu par
%touche Esc (touche d'echappement)
\begin{picture}(14,12)(0, 2) 
\put(0,0){\framebox(14,12)[lb]{}}
\put(1.5,2){\makebox(0,0)[lb]{\tiny Esc}}
\end{picture}\hspace{0.3ex},
%touche §
\begin{picture}(12,12)(0, 2)
\put(0,0){\framebox(12,12)[lb]{}}
\put(1.5,2){\makebox(0,0)[lb]{\tiny \S}}
\end{picture}\hspace{0.3ex}
ou
%touche #
\mbox{\begin{picture}(12,12)(0, 2)
\put(0,0){\framebox(12,12)[lb]{}}
\put(1.5,2){\makebox(0,0)[lb]{\tiny \#}}
\end{picture}\hspace{0.3ex}).}}\\ \hline
\end{tabular}
\vfil
\noindent\begin{tabular}{@{}|p{1cm}p{11.9cm}|@{}} \hline
\multicolumn{2}{|l|}{Musique~:  }  \\ \hline \hline
\idt   metronome$(n)$ & \apxFontSize \hspace{3em} règle  le  \underline{métronome}  à  $n$  battements  par  seconde  
             ({\idt largo} $=$ 50, 
                  {\idt larghetto} $=$ 63,    {\idt adagio} $=$ 71,   {\idt andante} $=$ 92,   {\idt moderato} $=$ 114, 
                  {\idt allegro} $=$ 144, {\idt presto} $=$ 184, {\idt prestissimo} $=$ 204).  \\
\idt   nuance$(n)$ & \apxFontSize \hspace{1em} règle la \underline{nuance} $n$  ({\idt pianissimo} $=$ 1,  {\idt piano} $=$ 2,  
                 {\idt mezzo\_piano} $=$ 3, {\idt un\_poco\_piano} $=$ 4,  {\idt sotto\_voce} $=$ 5, {\idt mezza\_voce} $=$ 6, 
				 {\idt un\_poco\_forte} $=$ 7, {\idt mezzo\_forte} $=$ 8, {\idt forte} $=$ 9, {\idt fortissimo} $=$ 10). \\
\idt   note$(h, d)$ & \apxFontSize \hspace{1em}\hspace{1ex} joue une \underline{note} de hauteur $h = (((o \ast octave) + n) + a)$ où $o = -3$, $-2$, 
               ..., $3$  est l'octave de la note ({\idt octave} $=$ 12),  $n$ $=$ {\idt d0} ou {\idt ut} $=$ 0, {\idt re} $=$ 2, {\idt mi} $=$ 4, {\idt fa} $=$ 5, 
			   {\idt sol} $=$ 7, {\idt la} $=$ 9, 
               {\idt si } $=$ 11 est le nom de la note,  $a$  =  {\idt double\_bemol} $= -2$,  {\idt bemol} $= -1$,  {\idt becarre} $= 0$,  
			   {\idt diese} $= 1$,  {\idt double\_diese} $= 2$ est l'altération de la note et  $d$ $=$ 
			   {\idt note\_carree} $= 8.0$, {\idt ronde} $= 4.0$, 
               {\idt ronde\_pointee} $= 6.0$, {\idt blanche} $= 2.0$, {\idt blanche\_pointee} $= 3.0$, 
			   {\idt blanche\_en\_triolet} $= 4/3$,  {\idt noire} $= 1.0$, 
               {\idt noire\_pointee} $= 1.5$,       {\idt noire\_en\_triolet} $= 2/3$,       {\idt croche} $= 0.5$, 
			   {\idt croche\_pointee} $= 0.75$, 
               {\idt croche\_en\_triolet} $= 1/3$,       {\idt double\_croche} $= 0.25$,       {\idt double\_croche\_pointee} $= 0.375$, 
               {\idt double\_croche\_en\_triolet} $= 0.5/3$,    {\idt triple\_croche} $= 0.125$,   {\idt triple\_croche\_pointee} $= 0.1875$, 
               {\idt triple\_croche\_en\_triolet} $= 0.25/3$,             {\idt quadruple\_croche} $= 0.0625$,           ou 
               {\idt quadruple\_croche\_en\_triolet} $= 0.125/3$ est la durée de la note.   \\
\idt   silence$(d)$ & \apxFontSize \hspace{2em} marque  un  \underline{silence}  de  durée   $d$  $=$
              {\idt baton\_de\_deux\_pauses} $= 8.0$,  {\idt pause} $= 4.0$, 
               {\idt pause\_pointee} $= 6.0$, {\idt demi\_pause\_pointee} $= 3.0$, {\idt demi\_pause} $= 2.0$, 
			   {\idt soupir} $= 1.0$, {\idt soupir\_pointe} $= 1.5$,
               {\idt demi\_soupir} $= 0.5$,         {\idt demi\_soupir\_pointe} $= 0.75$,            {\idt quart\_de\_soupir} $= 0.25$, 
              {\idt quart\_de\_soupir\_pointe} $= 0.375$, {\idt huitieme\_de\_soupir} $= 0.125$, 
			  {\idt huitieme\_de\_soupir\_pointe} $= 0.1875$,
              {\idt seizieme\_de\_soupir} $= 0.0625$.  \\ \hline
\end{tabular}
\vfil
\noindent\begin{tabular}{@{}|p{1cm}p{11.9cm}|@{}} \hline
\multicolumn{2}{|l|}{Interaction avec l'utilisateur du {\normalsize \rm\robot}~:  }  \\ \hline \hline
%\noindent\begin{tabular}{@{}|p{13.33cm}|@{}} \hline
%Interaction avec l'utilisateur du {\normalsize \rm\robot}~:  \\ \hline \hline
\multicolumn{2}{|l|}{\idt interaction$($X$,$ Y$,$ Ch$,$ Coordonnees$,$ Pivotements$,$ 
Dessin$,$ FinSouris$,$}\\
\idt Sauver$)$ & \apxFontSize \hspace{1em} retourne  les  coordonnées   $(x, y)$ du  {\normalsize \rm\robot}\   déplacé  par 
l'utilisateur avec la souris ou les flèches de  défilement
%fleche de defilement haut
\hspace{0.3ex}\begin{picture}(12,12)(0, 2)
\put(0,0){\framebox(12,12)[l]{$\scriptstyle\uparrow$}}
\end{picture}\hspace{0.3ex},
%fleche de defilement bas
\begin{picture}(12,12)(0, 2)
\put(0,0){\framebox(12,12)[l]{$\scriptstyle\downarrow$}}
\end{picture}\hspace{0.3ex},
%fleche de defilement gauche
\begin{picture}(12,12)(0, 2)
\put(0,0){\framebox(12,12)[bl]{$\scriptstyle\leftarrow$}}
\end{picture}\hspace{0.3ex}
ou
%fleche de defilement droite
\begin{picture}(12,12)(0, 2)
\put(0,0){\framebox(12,12)[bl]{$\scriptstyle\rightarrow$}}
\end{picture}\hspace{0.3ex} au moment où l'interaction avec l'utilisateur se termine. 
\\ \hline \end{tabular}
\clearpage
\noindent\begin{tabular}{@{}|p{1cm}p{11.9cm}|@{}} \hline
%\noindent\begin{tabular}{@{}|p{13.33cm}|@{}} \hline
%
& \apxFontSize L'interaction se termine quand
 le  bouton  de  la souris est enfoncé (auquel cas la variable {\idt Ch} de type {\idt char} prend pour valeur le
caractère nul {\idt chr$($0$)$}) ou quand une touche est tapée (le caractère correspondant étant affecté à {\idt Ch}). 
Les coordonnées $(x, y)$ du {\normalsize \rm\robot}\ sont relatives à sa position  au  moment 
           de l'appel de la procédure et comptées sur la grille.  Elles sont affectées aux variables {\idt X} et {\idt Y} de type
		   {\idt integer}.  L'interaction peut
être interrompue par les touches
%touche Esc (touche d'echappement)
\begin{picture}(14,12)(0, 2) 
\put(0,0){\framebox(14,12)[lb]{}}
\put(1.5,2){\makebox(0,0)[lb]{\tiny Esc}}
\end{picture}\hspace{0.3ex},
%touche §
\begin{picture}(12,12)(0, 2)
\put(0,0){\framebox(12,12)[lb]{}}
\put(1.5,2){\makebox(0,0)[lb]{\tiny \S}}
\end{picture}\hspace{0.3ex}
ou
%touche #
\begin{picture}(12,12)(0, 2)
\put(0,0){\framebox(12,12)[lb]{}}
\put(1.5,2){\makebox(0,0)[lb]{\tiny \#}}
\end{picture}\hspace{0.3ex}.
%

\tiret Si la valeur booléenne {\idt Coordonnees} est {\idt true} (respectivement {\idt false}) 
			alors les  coordonnées  courantes du {\normalsize \rm\robot}\ déplacé par la souris ou les flèches de défilement sont (ne sont pas) 
			affichées à l'écran pendant l'interaction.

\tiret Si la valeur booléenne {\idt Pivotements} est {\idt true} alors
           l'utilisateur a la  possibilité de faire pivoter interactivement le {\normalsize \rm\robot}\ (touches
%touche -
\begin{picture}(12,12)(0, 2)
\put(0,0){\framebox(12,12)[lb]{}}
\put(1.5,2){\makebox(0,0)[lb]{\tiny $-$}}
\end{picture}\hspace{0.3ex} ({\idt phtg}),
%touche G
\begin{picture}(12,12)(0, 2)
\put(0,0){\framebox(12,12)[lb]{}}
\put(1.5,2){\makebox(0,0)[lb]{\tiny G}}
\end{picture}\hspace{0.3ex} ({\idt pqtg}),
%touche +
\begin{picture}(12,12)(0, 2)
\put(0,0){\framebox(12,12)[lb]{}}
\put(1.5,2){\makebox(0,0)[lb]{\tiny $+$}}
\end{picture}\hspace{0.3ex}  ({\idt phtd}) et
%touche D
\begin{picture}(12,12)(0, 2)
\put(0,0){\framebox(12,12)[lb]{}}
\put(1.5,2){\makebox(0,0)[lb]{\tiny D}}
\end{picture}\hspace{0.3ex} ({\idt pqtd})) 
           avant la fin de l'interaction.  Dans ce cas l'interaction peut se terminer u\-ni\-que\-ment en enfonçant le bouton de la souris ou
		   la touche
%retour a la ligne
\framebox[11.mm][l]{$\hookleftarrow$~~~}\nolinebreak\hspace{0.3ex}\nolinebreak.

\tiret Si les valeurs booléennes {\idt Pivotements} et {\idt Dessin} sont égales à {\idt true} alors
	           l'utilisateur a la possibilité de dessiner avant la fin de l'interaction (touches
%touche A
\begin{picture}(12,12)(0, 2)
\put(0,0){\framebox(12,12)[lb]{}}
\put(1.5,2){\makebox(0,0)[lb]{\tiny A}}
\end{picture}\hspace{0.3ex} ({\idt av}),
%touche (
\begin{picture}(12,12)(0, 2)
\put(0,0){\framebox(12,12)[lb]{}}
\put(1.5,2){\makebox(0,0)[lb]{\tiny $($}}
\end{picture}\hspace{0.3ex} ({\idt vg}),
%touche )
\begin{picture}(12,12)(0, 2)
\put(0,0){\framebox(12,12)[lb]{}}
\put(1.5,2){\makebox(0,0)[lb]{\tiny $)$}}
\end{picture}\hspace{0.3ex} ({\idt vd}),
%touche )
\begin{picture}(12,12)(0, 2)
\put(0,0){\framebox(12,12)[lb]{}}
\put(1.5,2){\makebox(0,0)[lb]{\tiny $L$}}
\end{picture}\hspace{0.3ex} ({\idt lc}),
%touche )
\begin{picture}(12,12)(0, 2)
\put(0,0){\framebox(12,12)[lb]{}}
\put(1.5,2){\makebox(0,0)[lb]{\tiny $B$}}
\end{picture}\hspace{0.3ex} ({\idt bc}),
%touche )
\begin{picture}(12,12)(0, 2)
\put(0,0){\framebox(12,12)[lb]{}}
\put(1.5,2){\makebox(0,0)[lb]{\tiny $P$}}
\end{picture}\hspace{0.3ex} ({\idt dp}),
%touche )
\begin{picture}(12,12)(0, 2)
\put(0,0){\framebox(12,12)[lb]{}}
\put(1.5,2){\makebox(0,0)[lb]{\tiny $C$}}
\end{picture}\hspace{0.3ex} ({\idt dc}),
%touche )
\begin{picture}(12,12)(0, 2)
\put(0,0){\framebox(12,12)[lb]{}}
\put(1.5,2){\makebox(0,0)[lb]{\tiny $.$}}
\end{picture}\hspace{0.3ex}  ({\idt dg}) et
%touche )
\begin{picture}(12,12)(0, 2)
\put(0,0){\framebox(12,12)[lb]{}}
\put(1.5,2){\makebox(0,0)[lb]{\tiny $*$}}
\end{picture}\hspace{0.3ex}  (changement de la couleur du crayon) en corrigeant les erreurs avec
%touche d'effacement 
\framebox[11.mm][l]{$\leftarrow$~~~}\hspace{0.3ex}). Dans ce cas l'interaction peut se terminer en enfonçant le bouton de la souris ou
la touche
%retour a la ligne
\framebox[11.mm][l]{$\hookleftarrow$~~~}\nolinebreak\hspace{0.3ex}\nolinebreak.

\tiret Si la valeur booléenne {\idt FinSouris} est {\idt true} alors
	           l'utilisateur a la possibilité de terminer en enfonçant le bouton de la souris sinon
           il doit terminer en enfonçant une touche du clavier. 

\tiret Si les valeurs booléennes {\idt Dessin} et {\idt Sauver} sont {\idt true} alors un  programme  
		    {\idt Dessin\_du\_robot.p} dont l'exécution reproduit le dessin interactif est engendré sur disque.  S'il y a une erreur
(disquette vérrouillée,\ldots) alors {\idt  dessinDuRobotEngendre} sera égal à  {\idt false} après l'exécution de
{\idt interaction}.
\\ \hline
\end{tabular}
\newpage                 % Si on termine le chapitre par une page blanche
\thispagestyle{empty}
